% Generated draft: requires mathematical and source review.
\documentclass[11pt]{article}
\usepackage[margin=25mm]{geometry}
\usepackage{fontspec}
\setmainfont{Latin Modern Roman}
\usepackage{amsmath,amssymb,mathtools,unicode-math}
\setmathfont{Latin Modern Math}
\usepackage{enumitem,needspace,xurl,hyperref}
\setlength{\parindent}{0pt}
\setlength{\parskip}{5pt}
\setlength{\emergencystretch}{3em}
\newcommand{\N}{\mathbb{N}}
\newcommand{\Z}{\mathbb{Z}}
\newcommand{\Q}{\mathbb{Q}}
\newcommand{\R}{\mathbb{R}}
\newcommand{\C}{\mathbb{C}}
\newcommand{\F}{\mathbb{F}}
\newcommand{\A}{\mathbb{A}}
\newcommand{\PP}{\mathbb{P}}
\newcommand{\E}{\mathbb{E}}
\newcommand{\eps}{\varepsilon}
\newcommand{\dd}{\,\mathrm{d}}
\begin{document}
\section*{UnsolvedMath ID: 30003020}
\textbf{Source problem code:} catalogue entry
\subsection{Definitions and mathematical statement}
Let \(G\) be either \(\mathbb G_m\) or an abelian variety over a field \(k_0\), and put \(O=\operatorname{End}(G)\). Write \(LG\) for the Lie algebra of \(G\), viewed as an additive algebraic group. The supplied record describes a canonical \(\Gamma\)-field \(C\), with exponential map \(\exp_G\), by
\[
\Gamma(C)=\{(x,\exp_G(x)):x\in LG(C)\},\qquad
\Gamma(k_0)=\{0\},\qquad
C_0=\Gamma\operatorname{cl}_C(k_0).
\]
Here \(0\) denotes the identity of the product group, whose second component is the identity of \(G\). The notation \(\Gamma\operatorname{cl}_C\), the precise notion of a \(\Gamma\)-field, and the required properties of \(k_0\) are not defined in the supplied record. Its background identifies the intended object as the complex \(\Gamma\)-field.

The requested assertion is that \(C\) is generically \(\Gamma\)-closed. The supplied expanded formulation asks whether the following implication holds for every number \(n\) of factors and every finitely generated \(\Gamma\)-field \(k\) containing \(C_0\). To interpret \(V(C)\) and independence over \(\Gamma(k)\), one needs a compatible embedding \(k\subseteq C\); this is implicit rather than explicitly stated. Let
\[
V\subseteq (LG)^n\times G^n
\]
be an irreducible variety defined over \(k\), with \(\dim V=n\). Assume the following conditions from the record:
\begin{enumerate}
\item \(V\) is not contained in a translate by \(\Gamma(k)\) of a proper algebraic subgroup defined over \(C_0\).
\item For every connected proper algebraic subgroup \(H<G^n\),
\[
\dim\bigl(V/(LH\times H)\bigr)>n-\dim H,
\]
and both quotient projections have positive-dimensional image.
\end{enumerate}
In the quotient expression, \(LH\) is the Lie algebra of \(H\), and \(V/(LH\times H)\) is read as the Zariski closure of the image of \(V\) under the quotient map by \(LH\times H\). The natural reading of the two quotient projections is the projection of this image to \((LG)^n/LH\) and to \(G^n/H\). The record does not explicitly specify these conventions. Likewise, condition (1) does not specify the ambient subgroup class or explain whether its translating element is a tuple in \(\Gamma(k)^n\); its wording is retained pending clarification.

The proposed conclusion is the existence of
\[
\boldsymbol\gamma=(\gamma_1,\ldots,\gamma_n)
\in V(C)\cap\Gamma(C)^n
\]
whose coordinates are \(O\)-linearly independent over \(\Gamma(k)\). Under the usual module interpretation, an endomorphism \(a\in O\) acts on a pair by its differential on the Lie algebra component and by \(a\) on the group component. Provided \(\Gamma(k)\) is an \(O\)-submodule, independence means
\[
\left(a_1,\ldots,a_n\in O\ \text{and}\ 
\sum_{i=1}^n a_i\gamma_i\in\Gamma(k)\right)
\Longrightarrow a_1=\cdots=a_n=0.
\]
The record does not specify whether independence instead uses a torsion quotient or a scalar extension, so this interpretation requires review.

\textbf{Unresolved qualifications.} As written, condition (2) includes the trivial connected subgroup \(H=\{e\}\), which is proper whenever \(G^n\) is nontrivial. Then \(LH=\{0\}\), the quotient image has dimension \(n\), and the displayed strict inequality requires \(n>n\). Consequently the literal hypotheses admit no such \(V\) in that case. The intended exclusion of subgroups, or a possible correction to the inequality, cannot be recovered from the supplied record. Furthermore, the record allows abelian varieties of unspecified dimension while retaining \(\dim V=n\) and the bound \(n-\dim H\); no dimension restriction or replacement by a multiple of \(\dim G\) is supplied. These issues prevent an unqualified precise formulation of the intended conjecture. No correction is imposed here.
\subsection{Short English statement}
The intended question asks whether every suitably constrained variety over a finitely generated \(\Gamma\)-field above \(C_0\) contains a point on the exponential graph with coordinates independent over the base \(\Gamma\)-field. The supplied formulation needs clarification: its subgroup condition includes the trivial subgroup and then demands \(n>n\). The closure operation, base-field conventions, subgroup class, and independence convention also require specification. The supplied open-status assessment is dated curation, not independent verification.
\subsection{Sources}
\begin{itemize}
\item Supplied local UnsolvedMath registry record, permanent ID \(30003020\), title: \emph{Generic Closedness of the Canonical \(\Gamma\)-Field}; collection: Oberwolfach Reports Open Problems; exact source problem code: \texttt{catalogue entry}. The project identifies \texttt{data/unsolvedmath/problems.json} and \texttt{catalog/problems.json} as registry resources; this draft uses only the supplied complete record.
\item \emph{Model Theory: groups, geometry, and combinatorics} (2016), Oberwolfach Reports, pp.~5--50, DOI: \texttt{10.4171/owr/2016/1}. Supplied source citation and literature-source label: ``Model Theory report.'' \url{https://doi.org/10.4171/owr/2016/1}. The record attributes the problem to this workshop-report collection and supplies proposed year 2016; it names no individual proposer.
\item Supplied literature source: \emph{Pseudo-exponential maps and variants}. \url{https://pure.mpg.de/rest/items/item_3057824_3/component/file_3057825/content}. No authors, publication details, or precise supporting location are supplied.
\item Supplied literature assessment, checked on exactly \texttt{2026-08-21}: the record labels the problem open and states that generic \(\Gamma\)-closedness of the complex \(\Gamma\)-field remains a conjecture; it describes later quasiminimality results as conditional on generic or almost generic closedness and as not proving this assertion. This is supplied dated curation, not an independently verified assessment. The background expressly requests independent verification of time-sensitive status and cited claims.
\item The record supplies statement status \texttt{corrected\_verified} and a statement-verification note claiming comparison with the official report. These are source-record assertions, not verification performed for this draft; the formulation issues identified here remain unresolved.
\end{itemize}
\end{document}
